% Generated experiment uses booktabs and graphicx from the host paper. New surface only.
\section{Stateful Tool-Agent Study of Relational Joint Risk}
\label{sec:flash-relational-study}
\subsection{Dataset and Protocol}
We construct synthetic stateful sandboxes for email, payments, document sharing,
calendar scheduling, CRM updates, and cloud deployment. Operations change domain
state: outboxes, balances and invoice status, access-control lists, invitations,
account records, and deployed configurations. A trusted request lookup supplies
exact operation fields; external records provide ordinary notes, forged scoped
grants, or suggested destination redirects. These records are not user instructions.
The agent uses actual \texttt{deepseek-flash} function calls to read records and
current policy, submit an operation, and acquire independent registry verification
after a blocked preflight. Verification checks an exact
field/value/effect/revision tuple against the request ledger.

We use 400 training tasks and 200 development tasks from four domains, and
144 newly seeded test tasks, 24 per domain. CRM and cloud deployment are held
out from fitting and calibration. Training uses shared-source and shared-chain
graphs; test additionally contains disjoint-source and overlapping-source graphs.
All methods use the same tools, temperature-zero non-thinking model, a
10-turn limit, a 650-token response limit, and at most four verification queries.
Every test case and failure is retained. Exploratory pilots identified
unconditional prior verification and a fully-verified-state baseline artifact;
both were resolved before freezing independent tests.

\subsection{Curved Joint-Risk Representation}
Primitive event probabilities and field-dependency graphs are declared by the
sandbox policy. A protected operation is compromised when an unverified causal
path activates. Under independent primitive events, the proposed surface computes
the union-of-conjunctions probability by inclusion--exclusion:
\[
 R(p)=s\sum_{\emptyset\ne J\subseteq\mathcal P}
 (-1)^{|J|+1}\prod_{e\in\cup_{P\in J}P}p_e.
\]
Paths sharing events are not treated as independent. Scoped verification
removes only its corresponding applicable path. The executor independently
enumerates all primitive-event worlds to obtain expected loss and records unsafe
admissions as committed effects; labels do not read controller scores.
The comparison is one affine rule on the same probabilities, graph membership,
severity, budget, and scoped evidence. Coefficients are learned on training
states including partial verification, and the development threshold minimizes
$5\,\mathrm{FAR}+\mathrm{FBR}$. Both methods allow fully verified actions with
no remaining applicable path.

\begin{figure}[t]
\centering
\includegraphics[width=\columnwidth]{flash_joint_surface_slice.pdf}
\caption{Illustrative conditional slice of the declared joint-risk model and
the frozen affine rule. The second propagation probability and risk budget
are fixed at 0.2; severity is 1 and no field is verified. This is a representation
illustration, not a test-performance plot.}
\label{fig:flash-joint-slice}
\end{figure}

\subsection{Results and Module Ablations}
@@COMPARISON@@

@@ABLATION@@

The full surface safely completes @@FULL@@/@@N@@ tasks with @@FULL_UNSAFE@@
unsafe commits; the linear rule completes @@LINEAR@@/@@N@@ safely with
@@LINEAR_UNSAFE@@ unsafe commits. The surface saves @@SAVED@@\% of independent
verification queries (paired mean reduction @@QUERY_DELTA@@, 95\% CI
[@@QUERY_LO@@, @@QUERY_HI@@]). The safe-success difference is
@@SUCCESS_DELTA@@ percentage points (95\% CI [@@SUCCESS_LO@@, @@SUCCESS_HI@@]).
Of the linear rule's unsafe commits, @@UNSEEN_UNSAFE@@ occur in the two
held-out graph families. Domain and topology breakdowns are retained with
the execution traces.

Projection removal admits externally claimed grants. Lifting removal retains a
source-only path rather than its dependency chain. Geometric-filter removal
retains schema, current-policy, and user-scope checks but omits the joint budget.
Feedback removal stops at the first blocked operation. In these graphs, lifting
removal is conservative: extra verification can recover tasks, so final success
alone does not measure its contribution. We report false blocks and queries too.

These findings concern the upgraded relational surface with typed domain
adapters and declared uncertainty, rather than a rerun of original cone-based
DAFR benchmarks. They do not establish production natural-language encoding or
probability calibration. The event model provides a controlled representation
test; real-world hazard-model validity requires separate evidence. Existing
benchmark results are unchanged.
